\section{Monoidal categories}\label{sec:monoidal_categories}

\paragraph{Monoidal categories}

\begin{definition}\label{def:monoidal_category}\mcite[def. 6.1.1]{Perrone2024StartingCategoryTheory}
  A \term{monoidal category} is a \hyperref[def:category]{category} \( \cat{C} \) endowed with the following additional structure:
  \begin{thmenum}[series=def:monoidal_category]
    \thmitem{def:monoidal_category/product} Most importantly, we need a \hyperref[def:bifunctor]{bifunctor} \( {\otimes}: \cat{C}^2 \to \cat{C} \), which we will call (monoidal) \term{multiplication}.

    It is also called the \term{tensor product} because it generalizes \hyperref[def:module_tensor_product]{tensor products} of \hyperref[def:module]{modules}.

    Similarly to the \hyperref[def:hom_bifunctor]{\( \hom \)-bifunctor}, for every fixed object \( A \), by sending \( B \) to \( A \otimes B \) and \( g: B \to Y \) to \( \id_A \otimes g \), we obtain an endofunctor
    \begin{equation*}
      (A \otimes \anon*): \cat{C} \to \cat{C}.
    \end{equation*}

    Dually, by fixing the second argument to \( B \) (and \( \id_B \)), we obtain another endofunctor
    \begin{equation*}
      (\anon* \otimes B): \cat{C} \to \cat{C}.
    \end{equation*}

    Unlike with the \( \hom \)-functor, here both \( (A \otimes \anon*) \) and \( (\anon* \otimes B) \) are covariant.

    \thmitem{def:monoidal_category/unit} To emulate a \hyperref[def:semigroup_neutral_element]{neutral element}, we need a fixed object \( E \) of \( \cat{C} \), which we will call \term{unit}.

    \thmitem{def:monoidal_category/associator} In order to reconcile with the lack of proper \hyperref[def:binary_operation/associative]{associativity}, as in \cref{ex:def:natural_transformation/set_associator}, we need a \hyperref[def:natural_isomorphism]{natural isomorphism}, called the \term{associator}, with components
    \begin{equation*}
      \alpha_{A,B,C}: ((A \otimes B) \otimes C) \to (A \otimes (B \otimes C)).
    \end{equation*}

    The functors which \( \alpha_{A,B,C} \) transforms between are discussed in detail in \cref{ex:def:natural_transformation/set_associator} for the \hyperref[def:category_of_small_sets]{category of small sets}.

    The associator must satisfy the following pentagonal diagram:
    \begin{equation*}
      \begin{aligned}
        \includegraphics[page=1]{output/def__monoidal_category}
      \end{aligned}
    \end{equation*}

    Note that there are only two paths to check in this diagram, which are straightforward to check with a direct definition, like in \cref{ex:category_of_small_sets_is_monoidal}. With an indirect definition, like in \cref{thm:cartesian_monoidal_structure}, this check becomes difficult.

    \thmitem{def:monoidal_category/unitors} Finally, we need two more auxiliary natural isomorphisms, called the \term{left} and \term{right unitor}, with components
    \begin{align*}
      \lambda_A: (E \otimes A) \to A && \rho_A: (A \otimes E) \to A.
    \end{align*}

    The left unitor \( \lambda \) transforms \( (E \otimes \anon*) \) to the identity \( \id_{\cat{C}} \), while \( \rho \) transforms \( (\anon* \otimes E) \) to the identity.

    Their interaction with the associator is determined by the following triangular diagram:
    \begin{equation*}
      \begin{aligned}
        \includegraphics[page=2]{output/def__monoidal_category}
      \end{aligned}
    \end{equation*}
  \end{thmenum}
\end{definition}

\begin{example}\label{ex:category_of_small_sets_is_monoidal}
  We will show that the \hyperref[def:category_of_small_sets]{category of small sets} becomes \hyperref[def:monoidal_category]{monoidal} with multiplication given by the \hyperref[def:cartesian_product]{Cartesian product}. We will generalize this construction in \cref{thm:cartesian_monoidal_structure}, which will lead to the definition of Cartesian monoidal categories in \cref{def:cartesian_monoidal_category}.

  The product operation is given by \( A \times B \) for sets. To extend this definition for the functions \( {f: A \to X} \) and \( {g: B \to Y} \), we will use the definition
  \begin{equation*}
    \parens[\big]{ f \bincirc g }(a, b) \coloneqq \parens[\big]{ f(a), g(b) }.
  \end{equation*}

  Now consider the associator from \cref{ex:def:natural_transformation/set_associator}:
  \begin{equation*}
    \alpha_{A, B, C}\parens[\big]{ (a, b), c } \coloneqq \parens[\big]{ a, (b, c) }.
  \end{equation*}

  We have already shown that it is a natural isomorphism. Now we show that it satisfies the pentagonal diagram\fnote{This is not a formally a diagram since its vertices are members of the products, not the products themselves.}:
  \begin{equation*}
    \begin{aligned}
      \includegraphics[page=1]{output/ex__category_of_small_sets_is_monoidal}
    \end{aligned}
  \end{equation*}

  Consider the left unitor from \cref{ex:def:natural_transformation/set_left_unitor}, and the analogously defined right unitor:
  \begin{align*}
    \lambda_A(e, a) \coloneqq a,
    &&
    \rho_A(a, e) \coloneqq a.
  \end{align*}

  We have verified naturality, while invertibility is obvious, so these are natural isomorphisms. The triangular unitality diagram is easy to verify:
  \begin{equation*}
    \begin{aligned}
      \includegraphics[page=2]{output/ex__category_of_small_sets_is_monoidal}
    \end{aligned}
  \end{equation*}
\end{example}

\begin{definition}\label{def:strict_monoidal_category}\mcite[def. 6.1.3]{Perrone2024StartingCategoryTheory}
  We say that the \hyperref[def:monoidal_category]{monoidal category} is \term{strict} if its associator and unitors are identity transformations.

  We can thus omit them from the definition, and instead impose the following conditions:
  \begin{thmenum}
    \thmitem{def:strict_monoidal_category/unitality} We want the functors \( (E \otimes \anon*) \) and \( (\anon* \otimes E) \) to coincide with the identity \( \id_{\cat{C}} \):
    \begin{paracol}{2}
      \begin{leftcolumn}
        For every object \( A \), we want
        \begin{equation*}
          E \otimes A = A = A \otimes E.
        \end{equation*}
      \end{leftcolumn}

      \begin{rightcolumn}
        For every morphism \( g \), we want
        \begin{equation*}
          \id_E \otimes g = g = \id_E \otimes g.
        \end{equation*}
      \end{rightcolumn}
    \end{paracol}

    \thmitem{def:strict_monoidal_category/associativity} Strict associativity must hold for both objects and morphisms:
    \begin{paracol}{2}
      \begin{leftcolumn}
        For every triple of objects, we want
        \begin{equation*}
          (A \otimes B) \otimes C = A \otimes (B \otimes C).
        \end{equation*}
      \end{leftcolumn}

      \begin{rightcolumn}
        For every triple of morphisms, we want
        \begin{equation*}
          (f \otimes g) \otimes h = f \otimes (g \otimes h).
        \end{equation*}
      \end{rightcolumn}
    \end{paracol}
  \end{thmenum}
\end{definition}

\begin{example}\label{ex:def:strict_monoidal_category}
  We list some examples of \hyperref[def:strict_monoidal_category]{strict monoidal categories}:
  \begin{thmenum}
    \thmitem{ex:def:strict_monoidal_category/endofunctors} The most important example we will deal with are categories of endofunctors --- see \cref{def:endofunctor_category}.

    \thmitem{ex:def:strict_monoidal_category/monoid}\mcite[exerc. 6.1.8]{Perrone2024StartingCategoryTheory} Let \( \cat{M} \) be the \hyperref[def:discrete_category]{discrete category} induced by the underlying set \( U(M) \) of the \hyperref[def:monoid]{monoid} \( M \).

    With the unit \( e \) and with multiplication inherited from \( M \), \( \cat{M} \) becomes a strict monoidal category.

    Multiplication of objects with \( e \) is invariant in \( \cat{M} \) because it is invariant in \( M \), while for morphisms it is vacuous because \( \cat{M} \) only contains identity morphisms.

    Similarly, associativity of object multiplication in \( \cat{M} \) follows from that in \( M \), while for morphisms associativity is again vacuous.

    \thmitem{ex:def:strict_monoidal_category/monoid_order}\mcite[exerc. 6.1.9]{Perrone2024StartingCategoryTheory} A slightly more interesting situation occurs if \( M \) is an \hyperref[def:ordered_semigroup]{ordered monoid}, and we define \( \cat{M} \) as the \hyperref[def:thin_category]{thin category} of \( M \), as per \fullref{thm:order_category_isomorphism}. Compared to the previous example, \( \cat{M} \) has more morphisms --- namely, the morphism \( (x, y) \) in \( \cat{M} \) corresponds to the assumption \( x \leq y \) in \( M \).

    If \( a \leq b \) and \( x \leq y \), \cref{thm:def:ordered_semigroup/pairwise} implies that \( ax \leq by \). Thus,
    \begin{equation*}
      (a, b) \otimes (x, y) = (ax, by).
    \end{equation*}

    Multiplication with \( \id_e = (e, e) \) thus gives
    \begin{equation*}
      (e, e) \otimes (x, y)
      =
      (ex, ey)
      \reloset {\eqref{eq:def:semigroup_neutral_element/left}} =
      (x, y)
      \reloset {\eqref{eq:def:semigroup_neutral_element/right}} =
      (xe, ye)
      =
      (x, y) \otimes (e, e).
    \end{equation*}

    Furthermore, since multiplication in \( M \) is associative, it follows that
    \begin{align*}
      &\phantom{{}={}}
      \parens[\big]{ (x, y) \otimes (x', y') } \otimes (x^\dprime, y^\dprime)
      = \\ &=
      \parens[\big]{ (xx')x^\dprime, (yy')y^\dprime }
      \reloset {\eqref{eq:def:binary_operation/associative}} = \\ &=
      \parens[\big]{ x(x'x^\dprime), y(y'y^\dprime) }
      = \\ &=
      (x, y) \otimes \parens[\big]{ (x', y') \otimes (x^\dprime, y^\dprime) }.
    \end{align*}

    Therefore, the unitality and associativity conditions for \( \cat{M} \) are satisfied, and we conclude that \( \cat{M} \) is a strict monoidal category.
  \end{thmenum}
\end{example}

\begin{proposition}\label{thm:cartesian_monoidal_structure}\mcite[163]{MacLane1998CategoryTheory}
  Consider a \hyperref[def:category]{category} \( \cat{C} \) with \hyperref[def:categorical_product]{finite products}. We can use these to define a \hyperref[def:monoidal_category]{monoidal structure} on \( \cat{C} \).

  \begin{thmenum}
    \thmitem{thm:cartesian_monoidal_structure/product} For objects, we define \( A \otimes B \) as \( A \times B \). To extend this definition for the morphisms \( {f: A \to X} \) and \( {g: B \to Y} \), we will need to use the universal property of binary products from \cref{def:categorical_product/binary}:
    \begin{equation*}
      \begin{aligned}
        \includegraphics[page=1]{output/thm__cartesian_monoidal_structure}
      \end{aligned}
    \end{equation*}

    This defines a functor \( {\otimes}: \cat{C}^2 \to \cat{C} \).

    \thmitem{thm:cartesian_monoidal_structure/unitors} Fix a terminal object \( E \). We can use the right projection of \( E \times A \) as the left unitor \( \lambda_A \), and similarly the left projection of \( A \times E \) as the right unitor.

    \thmitem{thm:cartesian_monoidal_structure/associator} To obtain an associator \( \alpha \), we again use the universal property of products:
    \begin{equation*}
      \begin{aligned}
        \includegraphics[page=2]{output/thm__cartesian_monoidal_structure}
      \end{aligned}
    \end{equation*}
  \end{thmenum}
\end{proposition}
\begin{proof}
  We will skip proving the properties of \( \alpha \). Invertibility is simple enough, but naturality and the pentagon diagram require enormous diagrams (unlike the simple special case \( \cat{C} = \cat{Set} \) discussed in \cref{ex:category_of_small_sets_is_monoidal}). More abstract proof approaches are discussed in \cite{MathOF:proof_that_a_cartesian_category_is_monoidal}.

  \SubProof{Proof of functoriality of multiplication}

  \SubProofOf*{def:functor/CF1} By the universal property of \( A \times B \), the endomorphism \( \id_A \otimes \id_B \) is the unique one for which the following diagram commutes:
  \begin{equation*}
    \begin{aligned}
      \includegraphics[page=3]{output/thm__cartesian_monoidal_structure}
    \end{aligned}
  \end{equation*}

  But this diagram is also satisfied by \( \id_{A \times B} \), so the two are equal.

  \SubProofOf*{def:functor/CF2} We can stack the defining diagrams for \( f \otimes g \) and \( h \otimes k \):
  \begin{equation*}
    \begin{aligned}
      \includegraphics[page=4]{output/thm__cartesian_monoidal_structure}
    \end{aligned}
  \end{equation*}

  Incidentally, this also defines the morphism \( (h \bincirc f) \otimes (k \bincirc g) \). By the universal property of \( X \times Y \), the two coincide.

  \SubProof{Proof of naturality of the unitors} Naturality for \( f: A \to B \) follows from the defining diagram of \( \id_E \otimes f \):
  \begin{equation*}
    \begin{aligned}
      \includegraphics[page=5]{output/thm__cartesian_monoidal_structure}
    \end{aligned}
  \end{equation*}

  For \( \rho_A \), we proceed similarly.

  \SubProof{Proof of invertibility of the unitors} Denote the unique morphism from \( A \) to the terminal object \( E \) by \( e_A \). By the universal property of \( E \times A \), there exists a unique morphism \( \mu_A \) such that the following diagram commutes:
  \begin{equation*}
    \begin{aligned}
      \includegraphics[page=6]{output/thm__cartesian_monoidal_structure}
    \end{aligned}
  \end{equation*}

  This shows that \( \lambda_A \bincirc \mu_A = \id_A \), i.e. \( \mu_A \) is a right inverse of \( \lambda_A \).

  Since \( E \) is a terminal object, it has a unique map from \( E \times A \), so the projection \( \pi_E \) coincides with \( e_A \bincirc \lambda_A \). This allows us to extend the above diagram as follows:
  \begin{equation*}
    \begin{aligned}
      \includegraphics[page=7]{output/thm__cartesian_monoidal_structure}
    \end{aligned}
  \end{equation*}

  By the universal property of \( A \times B \), the composition \( \mu_A \bincirc \lambda_A \) is the unique morphism satisfying the above diagram. But it is also satisfied by \( \id_{E \times A} \), so the two coincide. It follows that \( \mu_A \) is a left inverse of \( \lambda_A \).

  By \cref{thm:def:morphism_invertibility/left_and_right}, \( \lambda_A \) is an isomorphism.

  For \( \rho_A \), we can proceed similarly.

  \SubProof{Proof of triangular diagram for unitors} We start with the following diagram, which glues together defining diagrams for \( \rho_A \otimes \id_B \) and \( \id_A \otimes \lambda_A \):
  \begin{equation*}
    \begin{aligned}
      \includegraphics[page=8]{output/thm__cartesian_monoidal_structure}
    \end{aligned}
  \end{equation*}

  To construct the associator \( \alpha_{A,E,B} \), we also need to project \( (A \times E) \times B \) to \( E \times B \):
  \begin{equation*}
    \begin{aligned}
      \includegraphics[page=9]{output/thm__cartesian_monoidal_structure}
    \end{aligned}
  \end{equation*}

  With the associator, we now have the unitality diagram:
  \begin{equation*}
    \begin{aligned}
      \includegraphics[page=10]{output/thm__cartesian_monoidal_structure}
    \end{aligned}
  \end{equation*}
\end{proof}

\paragraph{Endofunctor categories}

\begin{definition}\label{def:natural_transformation_horizontal_composition}\mcite[def. 1.4.24]{Perrone2024StartingCategoryTheory}
  Consider \hyperref[def:functor]{functors} and \hyperref[def:natural_transformation]{natural transformations} as in the diagram
  \begin{equation*}
    \begin{aligned}
      \includegraphics[page=1]{output/def__natural_transformation_horizontal_composition}
    \end{aligned}
  \end{equation*}

  We will construct a natural transformation \( (\beta \otimes \alpha): H \bincirc F \Rightarrow I \bincirc G \), called the \term{horizontal composition} of \( \alpha \) and \( \beta \).

  For this, note that, for any object \( A \) in \( \cat{C} \), the naturality of \( \beta \) gives us the following commutative square:
  \begin{equation*}
    \begin{aligned}
      \includegraphics[page=2]{output/def__natural_transformation_horizontal_composition}
    \end{aligned}
  \end{equation*}

  We simply let \( (\beta \otimes \alpha)_A \) be either of the two compositions above.
\end{definition}
\begin{comments}
  \item The notation {\( \otimes \)} highlights the role of horizontal composition as \hyperref[def:monoidal_category/product]{multiplication} in \hyperref[def:endofunctor_category]{categories of endofunctors}.
\end{comments}
\begin{defproof}
  The naturality diagram of \( \beta \otimes \alpha \) for any morphism \( f: A \to B \) can be constructed from those of \( H(\alpha) \) and \( \beta \):
  \begin{equation*}
    \begin{aligned}
      \includegraphics[page=3]{output/def__natural_transformation_horizontal_composition}
    \end{aligned}
  \end{equation*}
\end{defproof}

\begin{proposition}\label{thm:horizontal_composition_associativity}
  \hyperref[def:natural_transformation_horizontal_composition]{Horizontal composition} of natural transformations is associative in the following sense. Consider \hyperref[def:functor]{functors} and \hyperref[def:natural_transformation]{natural transformations} as in the diagram
  \begin{equation*}
    \begin{aligned}
      \includegraphics[page=1]{output/thm__horizontal_composition_associativity}
    \end{aligned}
  \end{equation*}

  Then the horizontal compositions of the following diagrams are equal:
  \begin{equation*}
    \includegraphics[page=2]{output/thm__horizontal_composition_associativity}
    \qquad\qquad
    \includegraphics[page=3]{output/thm__horizontal_composition_associativity}
  \end{equation*}

  Symbolically,
  \begin{equation*}
    (\gamma \otimes \beta) \otimes \alpha = \gamma \otimes (\beta \otimes \alpha).
  \end{equation*}
\end{proposition}
\begin{proof}
  We will use the path with the dashed arrow from the following diagram of \( (\beta \otimes \alpha)_A \):
  \begin{equation*}
    \includegraphics[page=4]{output/thm__horizontal_composition_associativity}
  \end{equation*}

  For \( J((\beta \otimes \alpha)_A) \), we have
  \begin{equation*}
    \includegraphics[page=5]{output/thm__horizontal_composition_associativity}
  \end{equation*}

  The dotted arrow is due to the naturality of \( \gamma \). With \( X = G(A) \), it coincides with the dotted arrow from the following diagram for \( (\gamma \otimes \beta)_X \):
  \begin{equation*}
    \includegraphics[page=6]{output/thm__horizontal_composition_associativity}
  \end{equation*}

  Therefore,
  \begin{equation*}
    [(\gamma \otimes \beta) \otimes \alpha]_A = [\gamma \otimes (\beta \otimes \alpha)]_A.
  \end{equation*}
\end{proof}

\begin{theorem}[Interchange law]\label{thm:composition_interchange_law}\mcite[prop. 1.4.25]{Perrone2024StartingCategoryTheory}
  \hyperref[def:functor_category/composition]{Vertical composition} \( {\bincirc} \) and \hyperref[def:natural_transformation_horizontal_composition]{horizontal composition} \( {\otimes} \) of \hyperref[def:natural_transformation]{natural transformations} can be interchanged.

  More precisely, consider natural transformations as in the diagram
  \begin{equation*}
    \begin{aligned}
      \includegraphics[page=1]{output/thm__composition_interchange_law}
    \end{aligned}
  \end{equation*}

  We claim that
  \begin{equation*}
    (\delta \bincirc \gamma) \otimes (\beta \bincirc \alpha) = (\delta \otimes \gamma) \bincirc (\beta \otimes \alpha).
  \end{equation*}
\end{theorem}
\begin{proof}
  From the naturality of \( \gamma \), we have
  \begin{equation*}
    \begin{aligned}
      \includegraphics[page=2]{output/thm__composition_interchange_law}
    \end{aligned}
  \end{equation*}

  Going through the top-right corner, we obtain the \( A \)-th component of \( (\delta \bincirc \gamma) \otimes (\beta \bincirc \alpha) \). Going through the bottom-left corner, we obtain the \( A \)-th component of \( (\delta \otimes \gamma) \bincirc (\beta \otimes \alpha) \).
\end{proof}

\begin{definition}\label{def:endofunctor_category}\mimprovised
  Unlike more general \hyperref[def:functor_category]{functor categories}, the category \( [\cat{C}, \cat{C}] \) of \hyperref[def:endofunctor]{endofunctors} on \( \cat{C} \) is a \hyperref[def:strict_monoidal_category]{strict monoidal category} with multiplication given by functor composition for objects and \hyperref[def:natural_transformation_horizontal_composition]{horizontal composition} for natural transformations.

  The multiplicative unit in \( [\cat{C}, \cat{C}] \) is, unsurprisingly, the identity functor.
\end{definition}
\begin{defproof}
  \SubProofOf{def:strict_monoidal_category/associativity} (Strict) associativity of object multiplication (functor composition) follows, as in \cref{def:category_of_small_categories}, from ordinary function composition associativity.

  Associativity of morphism multiplication (horizontal composition of natural transformations) is shown in \cref{thm:horizontal_composition_associativity}.

  \SubProofOf{def:strict_monoidal_category/unitality} Composing the endofunctor \( F \) with \( \id_{\cat{C}} \) from either side gives \( F \), so unitality for objects (functors) in \( [\cat{C}, \cat{C}] \) is obvious.

  To show left unitality of horizontal composition, consider functors and natural transformations as in the diagram
  \begin{equation*}
    \begin{aligned}
      \includegraphics[page=1]{output/def__endofunctor_category}
    \end{aligned}
  \end{equation*}

  The composition \( \id_{\id_{\cat{C}}} \otimes \alpha \) coincides with \( \alpha \) because its diagram at \( A \) is defined by
  \begin{equation*}
    \begin{aligned}
      \includegraphics[page=2]{output/def__endofunctor_category}
    \end{aligned}
  \end{equation*}

  Right unitality can be shown analogously.
\end{defproof}

\paragraph{Symmetric monoidal categories}

\begin{definition}\label{def:symmetric_monoidal_category}
  A \term{symmetric monoidal category} is a \hyperref[def:monoidal_category]{monoidal category} \( \cat{C} \) endowed with an additional \hyperref[def:natural_isomorphism]{natural isomorphism} with components
  \begin{equation*}
    \beta_{A,B}: A \otimes B \to B \otimes A.
  \end{equation*}

  We call \( \beta \) the \term{braiding} and require it to satisfy the equality
  \begin{equation*}
    \beta_{A,B}^{-1} = \beta_{B,A}
  \end{equation*}
  and the hexagonal identity
  \begin{equation*}
    \begin{aligned}
      \includegraphics[page=1]{output/def__symmetric_monoidal_category}
    \end{aligned}
  \end{equation*}
\end{definition}

\begin{example}\label{ex:category_of_small_sets_is_symmetric}
  Continuing \cref{ex:category_of_small_sets_is_monoidal}, we will show that the \hyperref[def:category_of_small_sets]{category of small sets} is \hyperref[def:symmetric_monoidal_category]{symmetric monoidal}. We will generalize this construction to Cartesian monoidal categories in \cref{thm:cartesian_symmetric_monoidal_structure}

  Consider the obvious braiding
  \begin{equation*}
    \beta_{A, B}(a, b) \coloneqq (b, a).
  \end{equation*}

  Clearly \( \beta_{A, B} \) is invertible with inverse \( \beta_{B, A} \).

  We can easily show naturality for functions \( f: A \to X \) and \( g: B \to Y \):
  \begin{equation*}
    \begin{aligned}
      \includegraphics[page=1]{output/ex__category_of_small_sets_is_symmetric}
    \end{aligned}
  \end{equation*}

  The braiding hexagon also commutes:
  \begin{equation*}
    \begin{aligned}
      \includegraphics[page=2]{output/ex__category_of_small_sets_is_symmetric}
    \end{aligned}
  \end{equation*}

  Therefore, \( \cat{Set} \) is symmetric monoidal.
\end{example}

\begin{proposition}\label{thm:cartesian_symmetric_monoidal_structure}
  Consider a \hyperref[def:category]{category} \( \cat{C} \) with \hyperref[def:categorical_product]{finite products}. In \cref{thm:cartesian_monoidal_structure}, we used these products to define a \hyperref[def:monoidal_category]{monoidal structure} on \( \cat{C} \).

  We can now use the universal property of \( A \times B \) to also define a braiding:
  \begin{equation*}
    \begin{aligned}
      \includegraphics[page=1]{output/thm__cartesian_symmetric_monoidal_structure}
    \end{aligned}
  \end{equation*}

  With this, \( \cat{C} \) becomes a \hyperref[def:symmetric_monoidal_category]{symmetric monoidal category}.
\end{proposition}
\begin{proof}
  We will skip proving the hexagonal identity of \( \beta \) because, similarly to the pentagonal diagram in \cref{thm:cartesian_monoidal_structure}, the required diagram is enormous.

  \SubProof{Proof that \( \beta_{A,B}^{-1} = \beta_{B,A} \)} We can stack the defining diagrams for \( \beta_{A, B} \) and \( \beta_{B, A} \):
  \begin{equation*}
    \begin{aligned}
      \includegraphics[page=2]{output/thm__cartesian_symmetric_monoidal_structure}
    \end{aligned}
  \end{equation*}

  By the uniqueness from universal property of \( A \times B \), the composition \( \beta_{B, A} \bincirc \beta_{A, B} \) coincides with the identity \( \id_{A \otimes B} \). Since the roles of \( A \) and \( B \) are symmetric, we conclude that \( \beta_{A, B} \bincirc \beta_{B, A} = \id_{B \otimes A} \).

  Therefore, \( \beta_{A, B} \) is fully invertible with inverse \( \beta_{B, A}^{-1} \).

  \SubProof{Proof of naturality of \( \beta \)} Fix morphisms \( f: A \to X \) and \( g: B \to Y \). We can stack the defining diagrams for \( \beta_{A, B} \) and \( \beta_{Y, X} \) with the defining diagram for \( g \otimes f \) from \cref{thm:cartesian_monoidal_structure/product}:
  \begin{equation*}
    \begin{aligned}
      \includegraphics[page=3]{output/thm__cartesian_symmetric_monoidal_structure}
    \end{aligned}
  \end{equation*}

  Incidentally, this is also the defining diagram for \( f \otimes g \), hence
  \begin{equation*}
    \begin{aligned}
      \includegraphics[page=4]{output/thm__cartesian_symmetric_monoidal_structure}
    \end{aligned}
  \end{equation*}

  This shows naturality of \( \beta \) for \( f \) and \( g \).
\end{proof}

\begin{definition}\label{def:cartesian_monoidal_category}\mcite[def. 6.1.13]{Perrone2024StartingCategoryTheory}
  We say that the \hyperref[def:monoidal_category]{monoidal category} is \term{Cartesian} if multiplication is given by \hyperref[def:categorical_product]{categorical products}, in accordance with \cref{thm:cartesian_monoidal_structure}.

  Cartesian categories are automatically equipped with the braiding transformation from \cref{thm:cartesian_symmetric_monoidal_structure}, so they are \hyperref[def:symmetric_monoidal_category]{symmetric monoidal}.
\end{definition}
\begin{comments}
  \item Every category with finite products has a Cartesian structure, but it may also have other monoidal structures.
\end{comments}

\paragraph{Closed monoidal categories}

\begin{definition}\label{def:closed_monoidal_category}\mcite[def. 6.5.1]{Perrone2024StartingCategoryTheory}
  A \term{closed monoidal category} is a \hyperref[def:symmetric_monoidal_category]{symmetric monoidal category} \( \cat{C} \) endowed with an additional \hyperref[def:functor]{functor}
  \begin{equation*}
    [\anon*, \anon*]: \cat{C}^\oppos \times \cat{C} \to \cat{C},
  \end{equation*}
  called the \term{internal \( \hom \)-functor}.

  For each object \( X \), we require \( [X, \anon*] \) to be \hyperref[def:adjoint_functor]{right adjoint} to the \hyperref[con:currying]{curried} product \( (\anon* \otimes X) \). If \( \cat{C} \) is \hyperref[def:small_category]{locally small}, and thus has well-defined \hyperref[def:hom_adjunction]{\( \hom \)-adjunctions}, so, for \( X = B \), the above condition amounts to the existence of a \hyperref[def:natural_isomorphism]{natural isomorphism} with components
  \begin{equation*}
    \varphi_{A, C}: \hom(A, [B, C]) \Rightarrow \hom(A \otimes B, C).
  \end{equation*}
\end{definition}

\begin{definition}\label{def:cartesian_closed_category}\mcite[394]{Perrone2024StartingCategoryTheory}
  If a \hyperref[def:cartesian_monoidal_category]{Cartesian monoidal category} \( \cat{C} \) has an \hyperref[def:closed_monoidal_category]{internal \( \hom \)-functor} \( [A, B] \), we call \( \cat{C} \) a \term{Cartesian closed category} and refer to \( [A, B] \) as an \term{exponential object} (alluding to the \hyperref[rem:set_of_functions_notation/power]{power notation} \( B^A \) for the set of all functions from \( A \) to \( B \)).
\end{definition}

\begin{proposition}\label{thm:category_of_small_sets_is_cartesian_closed}
  The \hyperref[def:category_of_small_sets]{category of small sets} is \hyperref[def:cartesian_closed_category]{Cartesian closed} with the \hyperref[def:set_of_all_functions]{set of all functions} \( \fun(A, B) = B^A \) acting as the internal \( \hom \)-functor.
\end{proposition}
\begin{proof}
  Due to \fullref{thm:cartesian_product_universal_property}, \( \cat{Set} \) has finite products, so, by \cref{thm:cartesian_symmetric_monoidal_structure}, it is symmetric monoidal.

  We describe an natural isomorphism from \( \fun(A \times B, C) \) to \( \fun(A, \fun(B, C)) \) in \cref{ex:def:natural_transformation/uncurrying_in_set}.

  Therefore, \( \cat{Set} \) is Cartesian closed.
\end{proof}

\begin{proposition}\label{thm:category_of_small_categories_is_cartesian_closed}
  The \hyperref[def:category_of_small_categories]{category of small categories} is \hyperref[def:cartesian_closed_category]{Cartesian closed} with the \hyperref[def:functor_category]{functor category} \( [\cat{A}, \cat{B}] \) acting as the internal \( \hom \)-functor.
\end{proposition}
\begin{proof}
  Due to \fullref{thm:product_category_universal_property}, \( \cat{Cat} \) has finite products, so, by \cref{thm:cartesian_symmetric_monoidal_structure}, it is symmetric monoidal.

  We describe an natural isomorphism from \( [\cat{A} \times \cat{B}, \cat{C}] \) to \( [\cat{A}, [\cat{B}, \cat{C}]] \) in \cref{ex:def:natural_transformation/uncurrying_in_cat}.

  Therefore, \( \cat{Cat} \) is Cartesian closed.
\end{proof}
